\chapter{Reference tables}
\label{app:tables}

\section{Common complexity classes}
\begin{center}
\begin{tabularx}{0.96\textwidth}{@{}p{0.16\textwidth}p{0.30\textwidth}X@{}}
\toprule
\textbf{Class} & \textbf{Typical example} & \textbf{Interpretation}\\
\midrule
$O(1)$ & hash lookup, under assumptions & bounded by input size in the model\\
$O(\log n)$ & binary search on sorted indexed data & repeated reduction of search space\\
$O(n)$ & linear scan & work proportional to input size\\
$O(n\log n)$ & comparison sorting families & near-linear times logarithmic growth\\
$O(n^2)$ & nested pair comparison & all or many pairs\\
\bottomrule
\end{tabularx}
\end{center}
These are growth descriptions, not benchmark results.

\section{Selected HTTP status meanings}
\begin{center}
\begin{tabular}{ll}
\toprule
200 & successful response\\
201 & resource created\\
204 & successful response with no content\\
400 & malformed or invalid request\\
401 & authentication required or failed\\
403 & authenticated but not authorised\\
404 & resource not found\\
409 & conflict with current state\\
422 & semantically invalid content, where used by the API\\
429 & rate limit exceeded\\
500 & unexpected server failure\\
503 & service temporarily unavailable\\
\bottomrule
\end{tabular}
\end{center}
The exact API contract should define which status is used for which condition.

\section{Evaluation-method selection}
\begin{center}
\begin{tabular}{p{0.25\textwidth}p{0.55\textwidth}}
\toprule
\textbf{Question} & \textbf{Useful method}\\
\midrule
Can experts identify interaction violations? & heuristic evaluation\\
Can a new user discover the task? & cognitive walkthrough or user study\\
What do people do in context? & observation or contextual inquiry\\
How do people describe trust or confusion? & interview or think-aloud\\
Does one version differ under a task? & comparative study with pre-defined outcomes\\
Does an organisational outcome change? & longitudinal or quasi-experimental design with process evidence\\
\bottomrule
\end{tabular}
\end{center}

\section{Security review prompts}
Identify assets, actors, trust boundaries, input validation, authorisation, secrets, logging, dependencies, rate limits, recovery, privacy, accessibility, and misuse. A checklist is a starting point; the threat model determines relevance.
